\subsection{基环的扩张：交换环情形}

现在设 A 是交换环，B 是环， \(\rho: \mathbf{A} \to \mathbf{B}\) 是环同态，使得 \(\rho(\mathbf{A})\) 包含于 B 的中心；换句话说， \(\rho\) 在 B 上定义一个

A-代数结构。对每个 A-模 E，右 B-模 \(\mathrm{E}_{(\mathbf{B})} = \mathrm{E} \otimes_{\mathbf{A}} \mathbf{B}\) 于是等同于 \(B \otimes_{A} E\) ，因为由假设， \(\rho_{*}(B_{s})\) 和 \(\rho_{*}(B_{d})\) 的 A-模结构相同。回顾对每对由 A-模组成的有序偶 (E, F)，我们定义了一个典范 B-同态

\[
\omega : \left(\operatorname{Hom} _ {\mathbf {A}} (\mathrm{E}, \mathrm{F})\right) _ {(\mathbf {B})} \rightarrow \operatorname{Hom} _ {\mathbf {B}} \left(\mathrm{E} _ {(\mathbf {B})}, \mathrm{F} _ {(\mathbf {B})}\right)\tag{10}
\]

它使得对所有 \(u \in \operatorname{Hom}_{\mathbf{A}}(\mathrm{E}, \mathrm{F})\) 有 \(\omega(u \otimes 1) = u \otimes 1_{\mathbf{B}}\) （《代数学》第二章 §5 第 3 目）。

命题 11. 设 A 是交换环，B 是环， \(\rho\) 是 A 到 B 的中心的同态，E 和 F 是两个 A-模。设 B 是平坦 A-模，且 E 是有限生成的（相应地有限表示的）。则典范同态 (10) 是单射的（相应地双射的）。

由于 \(\omega\) 是典范同构

\[
\operatorname{Hom} _ {\mathbf {A}} (\mathrm{E}, \mathbf {B} \otimes_ {\mathbf {A}} \mathrm{F}) \rightarrow \operatorname{Hom} _ {\mathbf {B}} (\mathrm{E} _ {(\mathbf {B})}, \mathrm{F} _ {(\mathbf {B})})
\]

与典范同态 (8)

\[
\nu \colon \mathbf {B} \otimes_ {\mathbf {A}} \operatorname{Hom} _ {\mathbf {A}} (\mathbf {E}, \mathbf {F}) \rightarrow \operatorname{Hom} _ {\mathbf {A}} (\mathbf {E}, \mathbf {B} \otimes_ {\mathbf {A}} \mathbf {F})
\]

（出处同上）的复合，该命题是第 9 目命题 10 的推论。

现在设 A 和 B 都是交换的，考虑三个 A-模 \(E_{1}\) 、 \(E_{2}\) 、 \(E_{3}\) 和一个 A-双线性映射 \(f: E_{1} \times E_{2} \to E_{3}\) 。则存在一个且仅一个 B-双线性映射 \(f_{B}: E_{1(B)} \times E_{2(B)} \to E_{3(B)}\) 使得

\[
f _ {\mathrm{B}} (1 \otimes x _ {1}, 1 \otimes x _ {2}) = 1 \otimes f (x _ {1}, x _ {2})
\]

对所有 \(x_{1} \in \mathbf{E}_{1}, x_{2} \in \mathbf{E}_{2}\) 成立（《代数学》第九章 §1 第 4 目命题 1）。

在下面的陈述中，我们将假设 B 是平坦 A-模，且对每个子模 \(E'\) （ \(E_{i}\) 的，i = 1, 2, 3），我们将典范地把 \(E'_{(B)}\) 与它在 \(E_{i(B)}\) 中的像等同（第 3 目，注记 2）。

命题 12. 设 A、B 是交换环， \(\rho\) 是 A 到 B 的同态，E \(_{1}\) 、 E \(_{2}\) 、 E \(_{3}\) 是三个 A-模， \(f\) : E \(_{1} \times\) E \(_{2} \to\) E \(_{3}\) 是一个 A-双线性映射，而

\[
f _ {\mathbf {B}} \colon \mathrm{E} _ {1 (\mathbf {B})} \times \mathrm{E} _ {2 (\mathbf {B})} \rightarrow \mathrm{E} _ {3 (\mathbf {B})}
\]

是它的扩张。考虑子模 \(\mathbf{F}_2\) （含于 \(\mathbf{E}_2\) ）与子模 \(\mathbf{F}_3\) （含于 \(\mathbf{E}_3\) ），并用 \(\mathbf{T}\) 表示 \(\mathbf{E}_1\) 中由那些 \(x_1 \in \mathbf{E}_1\) （它们满足 \(f(x_1, x_2) \in \mathbf{F}_3\) ，对所有 \(x_2 \in \mathbf{F}_2\) ）组成的子模。设 \(\mathbf{B}\) 是平坦 A-模，且 \(\mathbf{F}_2\) 是有限生成的。则 \(\mathbf{T}_{(\mathbf{B})}\) 是由那些 \(x_1' \in \mathbf{E}_{1(\mathbf{B})}\) （它们满足 \(f_{\mathbf{B}}(x_1', x_2') \in \mathbf{F}_{3(\mathbf{B})}\) ，对所有 \(x_2' \in \mathbf{F}_{2(\mathbf{B})}\) ）组成的集合。

设 \(p\) 是典范满射 \(\mathrm{E}_3 \to \mathrm{E}_3 / \mathrm{F}_3\) ；对每个 \(x_1 \in \mathrm{E}_1\) ，我们让它对应 A-线性映射 \(x_2 \mapsto p(f(x_1, x_2))\) （ \(\mathrm{F}_2\) 到 \(\mathrm{E}_3 / \mathrm{F}_3\) 的），记之为 \(g(x_{1})\) ；于是 g 是 \(E_{1}\) 到 \(\operatorname{Hom}_{\mathbf{A}}(\mathbf{F}_{2}, \mathbf{E}_{3}/\mathbf{F}_{3})\) 的一个 A-同态，且 g 的核恰好就是 T。由于 B 是平坦 A-模，我们有正合列

\[
0 \longrightarrow \mathrm{T} _ {(\mathbf {B})} \longrightarrow \mathrm{E} _ {1 (\mathbf {B})} \stackrel {1 \otimes g} {\longrightarrow} \left(\operatorname{Hom} _ {\mathbf {A}} \left(\mathrm{F} _ {2}, \mathrm{E} _ {3} / \mathrm{F} _ {3}\right)\right) _ {(\mathbf {B})}
\]

（第 3 目，命题 1）。由命题 11，典范同态

\[
\omega : \left(\operatorname{Hom} _ {\mathbf {A}} (\mathrm{F} _ {2}, \mathrm{E} _ {3} / \mathrm{F} _ {3})\right) _ {(\mathbf {B})} \rightarrow \operatorname{Hom} _ {\mathbf {B}} (\mathrm{F} _ {2 (\mathbf {B})}, \left(\mathrm{E} _ {3} / \mathrm{F} _ {3}\right) _ {(\mathbf {B})})
\]

是单射的。另一方面，由于 B 是平坦 A-模， \((\mathrm{E}_{3}/\mathrm{F}_{3})_{(\mathrm{B})}\) 典范地等同于 \(E_{3(\mathrm{B})}/F_{3(\mathrm{B})}\) ；取 \(\omega\) 与 \(1 \otimes g\) 的复合，我们得到一个同态 u，序列 \(^{1}\)

\[
0 \longrightarrow \mathrm{T} _ {(\mathrm{B})} \longrightarrow \mathrm{E} _ {1 (\mathrm{B})} \stackrel {u} {\longrightarrow} \operatorname{Hom} _ {\mathrm{B}} \left(\mathrm{F} _ {2 (\mathrm{B})}, \mathrm{E} _ {3 (\mathrm{B})} / \mathrm{F} _ {3 (\mathrm{B})}\right)
\]

是正合的。由定义立即可得， \(u(x_{1}^{\prime})\) （其中

\[
x _ {1} ^ {\prime} = 1 \otimes x _ {1} \in \mathrm{E} _ {1 (\mathbf {B})},
\]

）是这样的线性映射：它把每个 \(x_2' \in \mathrm{F}_{2(\mathbf{B})}\) 映到模 \(\mathrm{F}_{3(\mathbf{B})}\) 下的 \(f_{\mathbf{B}}(x_1', x_2')\) 的剩余类；由线性性，这对所有 \(x_1' \in \mathrm{E}_{1(\mathbf{B})}\) 也成立；由于 \(u\) 的核是 \(\mathrm{T}_{(\mathbf{B})}\) ，命题得证。

推论 1. 设 A、B 是两个交换环， \(\rho: \mathbf{A} \to \mathbf{B}\) 是同态，使得 B 是平坦 A-模，E 是有限表示 A-模。对 E 的每个有限生成子模 F， \(\mathrm{E}_{(\mathrm{B})}\) 的对偶中与 \(\mathrm{F}_{(\mathrm{B})}\) 正交的子模等于 \((\mathrm{F}')_{(\mathrm{B})}\) ，其中 \(\mathrm{F}'\) 是 E 的对偶 \(\mathrm{E}^*\) 中与 F 正交的子模。

由命题 11， \((\mathbf{E}^{*})_{(\mathbf{B})}\) 典范地同构于对偶 \((\mathbf{E}_{(\mathbf{B})})^{*}\) （即 \(\mathbf{E}_{(\mathbf{B})}\) 的对偶）。于是只需取 \(\mathbf{E}_1 = \mathbf{E}^*\) 、 \(\mathbf{E}_2 = \mathbf{E}\) 、 \(\mathbf{E}_3 = \mathbf{A}\) 、 \(\mathbf{F}_2 = \mathbf{F}\) 、 \(\mathbf{F}_3 = \{0\}\) ，并取 \(f\) 为 \(\mathbf{E}^{*} \times \mathbf{E}\) 上的典范双线性型，来应用命题 12。

推论 2. 设 A、B 是两个交换环， \(\rho: A \to B\) 是同态，使得 B 是平坦 A-模，E 是 A-模。则对 E 的每个有限生成子模 F， \(F_{(B)}\) 的零化理想是 B 的理想 aB，其中 a 是 F 在 A 中的零化理想。

只需取 \(\mathbf{E}_1 = \mathbf{A}\) 、 \(\mathbf{E}_2 = \mathbf{E}_3 = \mathbf{E}\) 、 \(\mathbf{F}_2 = \mathbf{F}\) 、 \(\mathbf{F}_3 = \{0\}\) 应用命题 12。

注记. 如果关于模 \(E_{t}\) 与双线性映射 f 都没有歧义，则有时用 \(F_{3}: F_{2}\) 表示命题 12 中记作 T 的模，并称它为 \(F_{2}\) 到 \(F_{3}\) 的迁移子。于是命题 12 的结论可读作

\[
\mathrm{F} _ {3 (\mathrm{B})}: \mathrm{F} _ {2 (\mathrm{B})} = \left(\mathrm{F} _ {3}: \mathrm{F} _ {2}\right) _ {(\mathrm{B})}.\tag{11}
\]

在诸 \(E_{i}\) 都等于环 A、f 是乘法且诸 \(F_{i}\) 为理想 \(a_{i}\) 的特殊情形，我们得到迁移子公式

\[
\mathbf {B} \left(a _ {3}: a _ {2}\right) = \mathbf {B} a _ {3}: \mathbf {B} a _ {2}\tag{12}
\]

它在 B 是平坦 A-模且 \(a_{2}\) 是有限生成理想时成立。
% semantic-fragments: legacy-input-seam
\relax{}
